\chapter*{Introduzione} %Se si cambia il Titolo cambiare anche la riga successiva così che appia corretto nell'indice
\addcontentsline{toc}{chapter}{Introduzione} %Per far apparire Introduzione nell'indice (Il nome deve rispecchiare quello del chapter)
\pagenumbering{arabic} % Settaggio numerazione normale

Nell'odierno panorama della cybersecurity, l'evoluzione costante delle minacce digitali richiede un approccio proattivo alla protezione delle informazioni e degli asset aziendali. In particolare, il \textbf{Dark Web}, la parte più profonda e pericolosa del deep web, accessibile solo tramite strumenti specializzati come TOR, rappresenta un ambiente in cui vengono spesso scambiate informazioni sensibili e dati sottratti in modo illecito, che possono rappresentare un rischio significativo per le organizzazioni. La disponibilità di dati e informazioni sul Dark Web costituisce una risorsa preziosa per le operazioni di \textbf{OSINT} (Open-source intelligence), che sfruttando fonti pubbliche e facilmente accessibili, permette di raccogliere informazioni e prevenire attacchi.\\

Il presente lavoro di tesi si inserisce in questo contesto, come strumento con l'obiettivo di esplorare l'integrazione delle informazioni del Dark Web come parte di una strategia di sicurezza aziendale, utilizzando risorse OSINT. Nello specifico, viene esaminato un caso di studio in cui una startup operante in ambito informatico e finanziario subisce un attacco. Questa minaccia porta alla compromissione delle credenziali, utilizzando strumenti di monitoraggio del Dark Web e di gestione dei log, viene dimostrato come un'organizzazione possa individuare le proprie vulnerabilità e adottare misure correttive e preventive.

Il contributo principale di questo elaborato risiede nella proposta di un approccio integrato che combina il Dark Web e l'uso di piattaforme OSINT, come soluzione per la gestione di incidenti e per rafforzare la sicurezza di un'azienda. Vengono inoltre esplorate le implicazioni pratiche all'adozione di tali strumenti, evidenziando come l'integrazione di diverse tecnologie possa migliorare la capacità di un'organizzazione di rilevare e rispondere alle minacce.\\

L'elaborato analizzerà come sia fondamentale mantenere una determinata preparazione e prontezza contro le minacce emergenti, adottando un approccio in costante evoluzione e incentrato sulla sicurezza. Questo include il monitoraggio regolare del Dark Web, la formazione costante dei dipendenti (per accrescere la consapevolezza sui rischi digitali) e l'aggiornamento delle policy di sicurezza. E verrà analizzata come l'integrazione di strumenti OSINT e di monitoraggio del Dark Web possano rappresentare un vantaggio competitivo significativo per le organizzazioni, aiutandole nella gestione dei rischi e nella protezione dei propri asset.
 

\afterpage{\null\newpage}